% ===============================
%          PREÁMBULO
% ===============================

% ---------- Idioma y codificación ----------
\usepackage[spanish,es-tabla]{babel}
\usepackage[T1]{fontenc}
\usepackage[utf8]{inputenc} % Si usas pdflatex
\usepackage{lmodern}        % Fuente estable (evita warnings de font shape)

% ---------- Márgenes ----------
\usepackage[
  left=3cm,
  right=2.5cm,
  top=3cm,
  bottom=3cm
]{geometry}

% ---------- Interlineado ----------
\usepackage{setspace}
\onehalfspacing   % Cambia a \doublespacing si lo exige la escuela

% ---------- Gráficos y tablas ----------
\usepackage{graphicx}
\graphicspath{{figures/}{logos/}}
\usepackage{booktabs}
\usepackage{longtable}
\usepackage{array}
\usepackage{caption}
\usepackage{subcaption}

% ---------- Matemáticas ----------
\usepackage{amsmath}
\usepackage{amssymb}

% ---------- Encabezados ----------
\usepackage{fancyhdr}
\pagestyle{fancy}
\fancyhf{}
\fancyhead[L]{Trabajo Terminal}
\fancyhead[R]{\leftmark}
\fancyfoot[C]{\thepage}
\setlength{\headheight}{16pt} % Elimina warning de fancyhdr

% ---------- Espaciado de capítulos ----------
\usepackage{titlesec}
\titlespacing*{\chapter}{-5pt}{0pt}{0pt}
% Si lo quieres aún más arriba: cambia 0pt por -5pt

% ---------- Hipervínculos ----------
\usepackage[hidelinks]{hyperref}
\usepackage{cleveref}

% ---------- Citas y bibliografía ----------
\usepackage{csquotes} % Necesario con babel + biblatex

\usepackage[
  backend=biber,
  style=ieee,
  sorting=none
]{biblatex}

\addbibresource{bib/referencias.bib}

% ---------- Mejor tipografía ----------
\usepackage{microtype}

% ===============================
%      FIN DEL PREÁMBULO
% ===============================